%% kapa: exam=1 round=24 year=2013 period=2 form=A subject=appraisal_law subject_ko=부동산관계법규 seq=19 number=59
%% source: Q-Net 감정평가사 기출문제 (cst00309) artlSeq=5073802 file=extracted/1008/5073802/2013년 제24회 감정평가사 1차시험 기출문제/2교시(A형).hwp
%% conversion: kordoc 4.24.51 -> tools/gen1.py (auto); review=proofread-sonnet; flags=none
\begin{question}{59}
부동산 가격공시 및 감정평가에 관한 법령상 감정평가사에 관한 설명으로 옳은 것은?
\begin{choices}
  \item 외국의 감정평가사 자격을 가진 자로서 「부동산 가격공시 및 감정평가에 관한 법률」에 의한 결격사유에 해당하지 아니하는 자는 그 본국에서 대한민국정부가 부여한 감정평가사 자격을 인정하지 않더라도 시ㆍ도지사의 인가를 받으면 감정평가업자의 업무를 행할 수 있다.
  \item 국유재산을 관리하는 기관에서 5년 이상 감정평가와 관련된 업무에 종사한 자로서 감정평가사 제2차시험에 합격한 자는 감정평가사의 자격이 있다.
  \item 감정평가사의 자격이 있는 자는 시ㆍ도지사에게 자격등록을 하면 감정평가업을 영위할 수 있다.
  \item 감정평가사는 감정평가업을 영위하기 위하여 복수의 사무소를 설치할 수 있다.
  \item 국토교통부장관은 감정평가사가 그 자격증을 다른 사람에게 양도한 경우 그 자격을 취소하여야 한다.
\end{choices}
\end{question}
\answer{2}
